% Options for packages loaded elsewhere
\PassOptionsToPackage{unicode}{hyperref}
\PassOptionsToPackage{hyphens}{url}
\PassOptionsToPackage{dvipsnames,svgnames,x11names}{xcolor}
\documentclass[
  american,
  10pt,
  a4paper,
]{article}
\usepackage{xcolor}
\usepackage[margin=24mm]{geometry}
\usepackage{amsmath,amssymb}
\setcounter{secnumdepth}{-\maxdimen} % remove section numbering
\usepackage{iftex}
\ifPDFTeX
  \usepackage[T1]{fontenc}
  \usepackage[utf8]{inputenc}
  \usepackage{textcomp} % provide euro and other symbols
\else % if luatex or xetex
  \usepackage{unicode-math} % this also loads fontspec
  \defaultfontfeatures{Scale=MatchLowercase}
  \defaultfontfeatures[\rmfamily]{Ligatures=TeX,Scale=1}
\fi
\usepackage{lmodern}
\ifPDFTeX\else
  % xetex/luatex font selection
\fi
% Use upquote if available, for straight quotes in verbatim environments
\IfFileExists{upquote.sty}{\usepackage{upquote}}{}
\IfFileExists{microtype.sty}{% use microtype if available
  \usepackage[]{microtype}
  \UseMicrotypeSet[protrusion]{basicmath} % disable protrusion for tt fonts
}{}
\makeatletter
\@ifundefined{KOMAClassName}{% if non-KOMA class
  \IfFileExists{parskip.sty}{%
    \usepackage{parskip}
  }{% else
    \setlength{\parindent}{0pt}
    \setlength{\parskip}{6pt plus 2pt minus 1pt}}
}{% if KOMA class
  \KOMAoptions{parskip=half}}
\makeatother
\usepackage{color}
\usepackage{fancyvrb}
\newcommand{\VerbBar}{|}
\newcommand{\VERB}{\Verb[commandchars=\\\{\}]}
\DefineVerbatimEnvironment{Highlighting}{Verbatim}{commandchars=\\\{\}}
% Add ',fontsize=\small' for more characters per line
\newenvironment{Shaded}{}{}
\newcommand{\AlertTok}[1]{\textcolor[rgb]{1.00,0.00,0.00}{\textbf{#1}}}
\newcommand{\AnnotationTok}[1]{\textcolor[rgb]{0.38,0.63,0.69}{\textbf{\textit{#1}}}}
\newcommand{\AttributeTok}[1]{\textcolor[rgb]{0.49,0.56,0.16}{#1}}
\newcommand{\BaseNTok}[1]{\textcolor[rgb]{0.25,0.63,0.44}{#1}}
\newcommand{\BuiltInTok}[1]{\textcolor[rgb]{0.00,0.50,0.00}{#1}}
\newcommand{\CharTok}[1]{\textcolor[rgb]{0.25,0.44,0.63}{#1}}
\newcommand{\CommentTok}[1]{\textcolor[rgb]{0.38,0.63,0.69}{\textit{#1}}}
\newcommand{\CommentVarTok}[1]{\textcolor[rgb]{0.38,0.63,0.69}{\textbf{\textit{#1}}}}
\newcommand{\ConstantTok}[1]{\textcolor[rgb]{0.53,0.00,0.00}{#1}}
\newcommand{\ControlFlowTok}[1]{\textcolor[rgb]{0.00,0.44,0.13}{\textbf{#1}}}
\newcommand{\DataTypeTok}[1]{\textcolor[rgb]{0.56,0.13,0.00}{#1}}
\newcommand{\DecValTok}[1]{\textcolor[rgb]{0.25,0.63,0.44}{#1}}
\newcommand{\DocumentationTok}[1]{\textcolor[rgb]{0.73,0.13,0.13}{\textit{#1}}}
\newcommand{\ErrorTok}[1]{\textcolor[rgb]{1.00,0.00,0.00}{\textbf{#1}}}
\newcommand{\ExtensionTok}[1]{#1}
\newcommand{\FloatTok}[1]{\textcolor[rgb]{0.25,0.63,0.44}{#1}}
\newcommand{\FunctionTok}[1]{\textcolor[rgb]{0.02,0.16,0.49}{#1}}
\newcommand{\ImportTok}[1]{\textcolor[rgb]{0.00,0.50,0.00}{\textbf{#1}}}
\newcommand{\InformationTok}[1]{\textcolor[rgb]{0.38,0.63,0.69}{\textbf{\textit{#1}}}}
\newcommand{\KeywordTok}[1]{\textcolor[rgb]{0.00,0.44,0.13}{\textbf{#1}}}
\newcommand{\NormalTok}[1]{#1}
\newcommand{\OperatorTok}[1]{\textcolor[rgb]{0.40,0.40,0.40}{#1}}
\newcommand{\OtherTok}[1]{\textcolor[rgb]{0.00,0.44,0.13}{#1}}
\newcommand{\PreprocessorTok}[1]{\textcolor[rgb]{0.74,0.48,0.00}{#1}}
\newcommand{\RegionMarkerTok}[1]{#1}
\newcommand{\SpecialCharTok}[1]{\textcolor[rgb]{0.25,0.44,0.63}{#1}}
\newcommand{\SpecialStringTok}[1]{\textcolor[rgb]{0.73,0.40,0.53}{#1}}
\newcommand{\StringTok}[1]{\textcolor[rgb]{0.25,0.44,0.63}{#1}}
\newcommand{\VariableTok}[1]{\textcolor[rgb]{0.10,0.09,0.49}{#1}}
\newcommand{\VerbatimStringTok}[1]{\textcolor[rgb]{0.25,0.44,0.63}{#1}}
\newcommand{\WarningTok}[1]{\textcolor[rgb]{0.38,0.63,0.69}{\textbf{\textit{#1}}}}
\usepackage{longtable,booktabs,array}
\usepackage{caption}
\captionsetup[table]{skip=6pt}
\newcounter{none} % for unnumbered tables
\usepackage{calc} % for calculating minipage widths
% Correct order of tables after \paragraph or \subparagraph
\usepackage{etoolbox}
\makeatletter
\patchcmd\longtable{\par}{\if@noskipsec\mbox{}\fi\par}{}{}
\makeatother
% Allow footnotes in longtable head/foot
\IfFileExists{footnotehyper.sty}{\usepackage{footnotehyper}}{\usepackage{footnote}}
\makesavenoteenv{longtable}
\ifLuaTeX
\usepackage[bidi=basic,shorthands=off]{babel}
\else
\usepackage[bidi=default,shorthands=off]{babel}
\fi
\ifLuaTeX
  \usepackage{selnolig} % disable illegal ligatures
\fi
\setlength{\emergencystretch}{3em} % prevent overfull lines
\providecommand{\tightlist}{%
  \setlength{\itemsep}{0pt}\setlength{\parskip}{0pt}}
\usepackage{microtype}
\usepackage{xurl}
\setlength{\emergencystretch}{3em}
\usepackage{bookmark}
\IfFileExists{xurl.sty}{\usepackage{xurl}}{} % add URL line breaks if available
\urlstyle{same}
% fallback for those not using the hyperref driver hyperxmp:
\makeatletter
\@ifundefined{xmpquote}{\newcommand{\xmpquote}[1]{#1}}{}
\makeatother
\hypersetup{
  pdftitle={Security Analysis of ARCANE-Octarine},
  pdfauthor={Mounir IDRASSI},
  pdflang={en-US},
  colorlinks=true,
  linkcolor={MidnightBlue},
  filecolor={Maroon},
  citecolor={Blue},
  urlcolor={MidnightBlue},
  pdfcreator={LaTeX via pandoc}}

\title{Security Analysis of ARCANE-Octarine}
\usepackage{etoolbox}
\makeatletter
\providecommand{\subtitle}[1]{% add subtitle to \maketitle
  \apptocmd{\@title}{\par {\large #1 \par}}{}{}
}
\makeatother
\subtitle{Concrete attack bounds, proof obligations, and specification
consistency --- Version 1.0.1}
\author{Mounir IDRASSI}
\date{23 September 2026}

\begin{document}
\maketitle

\section*{Abstract}\label{abstract}
\addcontentsline{toc}{section}{Abstract}

ARCANE-Octarine is a Fiat--Shamir signature scheme using learning with
rounding over a radical ring and a power-of-two modulus. This document
analyzes the submitted 128-, 256-, and 512-level profiles, their
concrete hash instantiations, the relations used in the security
estimates, and the supplied implementations. It gives a chosen-message
signature-transfer attack through the SM3 counter-mode hash
construction, a secret-free challenge-search forgery, a 344-bit
description of an equivalent signing secret in the SM3 family, and a
polynomial-time solution of the relaxed homogeneous SIS instances used
in the 512-level estimates. The last result is an attack on the
estimated auxiliary problem; its witness violates the stricter response
bound of the actual signature relation. The document also identifies
missing quantitative proof arguments, modular structure requiring
analysis, and deterministic disagreements between Algorithm 2 and all
implementation families at levels 128 and 256.

The evidence does not establish a practical forgery or recovery of an
unknown signing key at the submitted parameters. It does establish that
the advertised security claims are unsupported, with explicit generic
attack bounds far below the 256- and 512-bit classical targets for the
SM3 implementations. Neither supplied hash family supports 512-bit
message-collision strength. Continued consideration requires revised
primitives, parameters, normative specifications, and analysis of the
actual verification relations.

\clearpage
\begingroup
\small
\setlength{\parskip}{0pt}
\setcounter{tocdepth}{2}
\tableofcontents
\endgroup
\clearpage

\section{1. Scope, notation, and
interpretation}\label{scope-notation-and-interpretation}

\subsection{1.1 Evaluated artifact}\label{evaluated-artifact}

The subject is \emph{ARCANE-Octarine: Algorithm Specifications and
Supporting Documentation}, attributed in the submission to the ARCANE
team {[}S{]}, together with the supplied reference SM3, optimized SM3,
and additional SHA3/SHAKE implementation families. The specification is
a 53-page PDF, identified by SHA-256:

\begin{Shaded}
\begin{Highlighting}[]
\NormalTok{b88fbdd019beb8f0ae85c1554b48c87a42c5447b39171ce559ae30988c1bf09f}
\end{Highlighting}
\end{Shaded}

All specification page references below use printed page numbers; the
corresponding one-based PDF page number is two greater. Findings apply
to this snapshot and the source files identified in the evidence
manifest {[}E9{]}. They do not presume that a future revision has the
same behavior.

Author contact:
\href{mailto:mounir@amcrypto.jp}{\nolinkurl{mounir@amcrypto.jp}}. Report
and original results: CC BY 4.0; research software: MIT. The
accompanying repository is
\href{https://github.com/amcrypto-jp/octarine-cryptanalysis}{octarine-cryptanalysis}.
The evidence inventory in Section 12 identifies the programs and data
needed to reproduce the checks.

An \emph{exact construction} below is a mathematical algorithm with
stated hypotheses. A \emph{generic work bound} assumes the indicated
idealized hash or search model. An \emph{estimator output} describes a
modeled attack, and does not lower-bound the cost of every possible
attack. A \emph{diagnostic} tests one property or example and does not
establish security.

\subsection{1.2 Parameters and
verification}\label{parameters-and-verification}

Write

\[
R_q=(\mathbb Z/q\mathbb Z)[y]/(y^{1024}+1),\qquad
S_q=R_q[x]/(x^k-y-2),\qquad N=1024k,
\]

where \(q=2^{24}\). Norms refer to coefficient infinity norms in the
submitted \((y^i x^j)\) representation. Put \(g=q/p\) and
\(m=2\gamma_2\). A challenge \(c\) is a polynomial in \(R_q\), embedded
in the constant \(x\) component of \(S_q\).

{\def\LTcaptype{none} % do not increment counter
\begin{longtable}[]{@{}lrrr@{}}
\toprule\noalign{}
Parameter & Octarine-128 & Octarine-256 & Octarine-512 \\
\midrule\noalign{}
\endhead
\bottomrule\noalign{}
\endlastfoot
\(k\) & 1 & 2 & 5 \\
\(N\) & 1024 & 2048 & 5120 \\
\(p\) & \(2^{21}\) & \(2^{22}\) & \(2^{22}\) \\
\(\eta=q/(2p)\) & 4 & 2 & 2 \\
\(d\) & 11 & 13 & 14 \\
\(\tau\) & 16 & 36 & 87 \\
\(\gamma_1\) & \(2^{18}\) & \(2^{19}\) & \(2^{21}\) \\
\(\gamma_2\) & \(2^{17}\) & \(2^{18}\) & \(2^{20}\) \\
\(\beta=\tau\eta\) & 64 & 72 & 174 \\
\(\omega\) & 79 & 210 & 399 \\
Challenge-seed bytes & 32 & 64 & 128 \\
Claimed classical / quantum bits & 128 / 80 & 256 / 128 & 512 / 256 \\
\end{longtable}
}

Key generation produces \(s_1\) with coefficients in \([-\eta,\eta-1]\),
and defines

\[
b=\lfloor(p/q)As_1\rceil,\qquad
s_2=As_1-gb,\qquad b=b_1 2^d+b_0.
\]

The public key contains the seed for \(A\) and the compressed value
\(b_1\). For a message \(M\), the verifier computes

\[
tr=H_{tr}(pk),\quad \mu=H_\mu(tr,M),\quad
c=\operatorname{SampleInBall}(\widetilde c),
\]

\[
v=Az-cgb_1 2^d,\qquad
w_1=\operatorname{UseHint}(h,v,m).
\]

After canonical decoding, acceptance requires

\[
\|z\|_\infty<\gamma_1-\beta,\qquad
\operatorname{HW}(h)\leq\omega,\qquad
\widetilde c=H_{ch}(\mu,w_1).
\]

Domain labels, output lengths, and length-delimited fields are implicit
in these hash expressions. All three profiles have 128-byte \(tr\) and
\(\mu\) values. The stored challenge seed has \(2\lambda\) bits, where
\(\lambda\) is the profile label.

EUF-CMA means existential unforgeability under chosen-message attack;
SUF-CMA additionally prohibits a fresh signature on a previously signed
message. Every EUF-CMA attack is also a SUF-CMA attack.

\section{2. Main findings and their
scope}\label{main-findings-and-their-scope}

{\def\LTcaptype{none} % do not increment counter
\begin{longtable}[]{@{}
  >{\raggedright\arraybackslash}p{(\linewidth - 4\tabcolsep) * \real{0.3333}}
  >{\raggedright\arraybackslash}p{(\linewidth - 4\tabcolsep) * \real{0.3333}}
  >{\raggedright\arraybackslash}p{(\linewidth - 4\tabcolsep) * \real{0.3333}}@{}}
\toprule\noalign{}
\begin{minipage}[b]{\linewidth}\raggedright
Finding
\end{minipage} & \begin{minipage}[b]{\linewidth}\raggedright
Established result
\end{minipage} & \begin{minipage}[b]{\linewidth}\raggedright
Scope limit
\end{minipage} \\
\midrule\noalign{}
\endhead
\bottomrule\noalign{}
\endlastfoot
SM3 message-hash construction & An equal-length internal-state collision
transfers a signature; generic classical work is about \(2^{128}\). A
48-bit reduced-state experiment executes the transfer end to end & The
experiment changes the SM3 primitive in every hash role; it is a
reduced-state model, not a forgery against the submitted 256-bit
profile \\
Challenge search & A fixed challenge with \(z=h=0\) gives a complete
generic forgery strategy & Quantum figures count oracle queries, not
logical gates \\
SM3 secret expansion & A 344-bit descriptor generates the complete
\(s_1\); equivalent-key search needs at most \(2^{344}\) candidates & No
unknown-key search was performed \\
512-level relaxed SIS & A nonzero norm-\(q/4\) vector is found by binary
linear algebra with constant success probability & The vector fails the
actual stricter response bound \\
Proof instantiation & Statistical terms, simulator behavior, seeded
sampling, and relation modeling are not concretely discharged & This
does not prove that a suitable proof is impossible \\
Ring structure & Nonunit mass, a two-component decomposition at \(k=5\),
and even-weight challenge restrictions are explicit & A useful
projection or transcript attack remains unestablished \\
Specification conformance & Algorithm 2 and every source family use
different challenge maps at levels 128 and 256 & This alone does not
show lower challenge entropy \\
\end{longtable}
}

The central recommendation is to withhold endorsement of the submitted
security levels and withdraw the present 512-bit claim. Sections 3--6
give the attack constructions and bounds. Sections 7--10 distinguish the
remaining proof, algebraic, and implementation issues from attacks.

\section{3. Hash instantiation and chosen-message signature
transfer}\label{hash-instantiation-and-chosen-message-signature-transfer}

\subsection{3.1 The implemented counter
construction}\label{the-implemented-counter-construction}

The specification names an SM3-based pseudoXOF for the reference and
optimized families, but does not give its full algorithm on pp.~16--17.
The supplied function \emph{pseudoXOF} in the corresponding
\emph{utils/auxfunc.c} files, starting at line 482, implements

\[
\operatorname{XOF}_{\ell}(X)=
\operatorname{prefix}_{\ell}\big(
\operatorname{SM3}(X\Vert\operatorname{ctr}_{32}(1))
\Vert\operatorname{SM3}(X\Vert\operatorname{ctr}_{32}(2))
\Vert\cdots\big),
\]

where \(\ell\) is the output length in bits and the counter is encoded
big-endian. The parameter macros bind this construction to the signature
hash roles and to public, secret, mask, and challenge expansion.

SM3 uses a 256-bit chaining value and 512-bit message blocks. Suppose
two distinct equal-length strings reach the same chaining value at the
same block boundary, before a common suffix. For every counter value,
the subsequent input blocks and length padding are identical.
Determinism then implies equality of every counter output block, and
hence of the entire requested XOF output.

This argument requires a collision before the common suffix. Equality of
only the first final digest block, without equality of the relevant
internal state, would not suffice.

For an idealized 256-bit state map, birthday search gives about
\(2^{128}\) classical evaluations. Generic quantum collision algorithms
give an approximately \(2^{256/3}=2^{85.33}\) query bound, with
substantial storage and coherent-computation requirements {[}BHT98{]}.
These are generic resource bounds, not measured attack implementations.
For an output shorter than 256 bits, its ordinary output-length
collision bound must also be included.

The toy program {[}E3{]} exhibits a state collision and its propagation
through 64 output blocks. The C program {[}E4{]} checks SM3 continuation
against the submitted implementation. It forces the same internal state
to validate propagation; it does not find colliding full-round SM3
messages.

A concrete reduced-step example replays the two second blocks from
Mendel, Nad, and Schläffer {[}MNS13, Table 3{]}. The blocks differ in
eight bytes. Starting from the table's supplied chaining input \(h_1\),
an adaptation of the submitted \emph{sm3\_bit\_compress} with 20 steps
and XOR feed-forward sends both blocks to

\texttt{b2033829677c16d2a6de9db9fd898668a9119d20476364d6a0838adc08d3833d}

matching the published \(h_2\). The adapter's 64-step outputs are
cross-checked against the unmodified submitted function, which gives
different states for this pair. The computed 20-step states retain
equality through common padding and 64 counter continuations in the
custom-IV variant {[}E4a{]}.

This replay supplies \(h_1\) directly; it does not establish a path to
\(h_1\) from the standard SM3 IV or from the state after Octarine's
public-key-dependent prefix. The submitted hash executes 64 steps, and
no Octarine signing or verification operation is run. The example is an
attributed illustration of collision propagation, not a new collision
attack or an improvement over the generic bound for Octarine.

A separate experiment now executes the transfer mechanism end to end in
a deliberately reduced-state model {[}E4b{]}. It starts from the
hash-verified Octarine-256 reference utility and changes the SM3
primitive in two ways: after each compression it retains only the first
48 chaining-state bits, and it serializes the live six state bytes by
repetition across the 32-byte digest. This serialization is injective on
the reduced state. Both changes affect every SM3 use in the experiment;
Octarine's domain encoders and key-generation, signing, and verification
routines are unchanged. The resulting hash is not the submitted SM3
profile.

For a deterministic test key, the search finds two distinct 64-byte
blocks that collide at block offset 3 of the actual \(H_\mu\) input
after 12,894,937 candidate blocks. The resulting equal-length 120-byte
messages have an identical 1024-bit \(\mu\). One signature query on the
first message yields a signature accepted for both it and the
never-signed second message. A different-suffix message, a truncated
signature, and an independently checked full-length corrupted signature
are rejected. The independent verifier consumes only the public key,
signature, and messages; the generated test secret key is not included
in the recorded campaign evidence.

At 256 bits the instrumentation's reduced-state branches are inactive:
known-answer tests and deterministic key, signature, and verification
outputs match the submitted utility. Preprocessed code also matches
after normalizing only source-file paths and line numbers embedded in
error diagnostics. Replaying the recorded pair with the submitted
256-bit SM3 instead gives different states and different \(\mu\) values;
the signature does not transfer in that full-width control. Birthday
trials at widths 8 through 48 are recorded as finite-width observations.
They illustrate search scaling but do not empirically establish a cost
at 256 bits; the approximately \(2^{128}\) estimate follows from the
generic birthday argument.

This model experiment supports the end-to-end connection between an
internal-state collision and signature transfer. It is not a full-width
collision, does not transfer a signature under the submitted parameters,
and does not reduce the generic work factor for those parameters. The
20-step published collision above is a separate illustration and should
not be confused with this scaled-state campaign.

\subsection{3.2 Consequence for
unforgeability}\label{consequence-for-unforgeability}

For a fixed public key, the signature uses \(M\) only through
\(\mu=H_\mu(tr,M)\). An adversary can choose sufficiently long,
equal-length messages with a common alignment prefix and vary a complete
message block. The encoded length fields remain identical. A collision
of the internal state after that block, followed by the same suffix,
makes the complete 1024-bit \(\mu\) equal.

The attack is:

\begin{enumerate}
\def\labelenumi{\arabic{enumi}.}
\tightlist
\item
  Find distinct messages \(M_0,M_1\) with the same \(\mu\) using the
  state-collision construction.
\item
  Request a signature \(\sigma\) on \(M_0\).
\item
  Output \((M_1,\sigma)\).
\end{enumerate}

The verifier performs identical computations for the two messages, so
\(\sigma\) verifies for the never-signed \(M_1\). The construction
breaks EUF-CMA and message binding at the collision-search cost. It
needs one signing query. Length delimiters and domain separation prevent
ambiguity between inputs; they do not prevent equal-length collisions
within one domain.

Consequently, the SM3 families have a generic classical
signature-transfer bound near \(2^{128}\) at every profile. This is at
the 128-level target, and far below the 256- and 512-level targets. It
does not by itself establish an attack in a two-key exclusive-ownership
game.

\subsection{3.3 SHAKE and per-role
requirements}\label{shake-and-per-role-requirements}

The additional implementation uses SHAKE256 for signature hashes and
secret/challenge expansion, and SHAKE128 for public expansion. FIPS 202,
Appendix A, gives SHAKE256 collision strength
\(\min(d_{\mathrm{out}}/2,256)\) for \(d_{\mathrm{out}}\) output bits
{[}F202{]}. Thus the 1024-bit message representative has generic
collision strength capped at 256 bits. Increasing its output length
cannot supply the claimed 512-bit classical message-binding strength.

The preimage entry in FIPS 202 is a lower bound. It should not be
restated as a universal 256-bit upper bound on every SHAKE256 preimage
problem. The conclusion here concerns collisions in the message-hash
role.

{\def\LTcaptype{none} % do not increment counter
\begin{longtable}[]{@{}
  >{\raggedright\arraybackslash}p{(\linewidth - 4\tabcolsep) * \real{0.3333}}
  >{\raggedleft\arraybackslash}p{(\linewidth - 4\tabcolsep) * \real{0.3333}}
  >{\raggedright\arraybackslash}p{(\linewidth - 4\tabcolsep) * \real{0.3333}}@{}}
\toprule\noalign{}
\begin{minipage}[b]{\linewidth}\raggedright
Backend for \(H_\mu\)
\end{minipage} & \begin{minipage}[b]{\linewidth}\raggedleft
Generic classical collision exponent
\end{minipage} & \begin{minipage}[b]{\linewidth}\raggedright
Consequence
\end{minipage} \\
\midrule\noalign{}
\endhead
\bottomrule\noalign{}
\endlastfoot
SM3 counter construction & 128 & Insufficient for 256- and 512-bit
claims \\
SHAKE256, 1024-bit output & 256 & Insufficient for a 512-bit claim \\
\end{longtable}
}

The specification's generic requirement of \(\lambda/2\) collision bits
and \(\lambda\) preimage bits is insufficient for this role. If
\(\lambda\) denotes a claimed classical unforgeability work factor, the
signature-transfer route must itself cost at least \(2^\lambda\) in the
chosen resource convention. A per-role security table is necessary; a
long output is not evidence that the underlying construction meets that
table.

\section{4. A secret-free challenge-search
forgery}\label{a-secret-free-challenge-search-forgery}

\subsection{4.1 Construction}\label{construction}

Let \(C_\tau\) be the set of signed weight-\(\tau\) polynomials in
\(R_q\):

\[
|C_\tau|=\binom{1024}{\tau}2^\tau.
\]

Choose any fixed \(c_0\in C_\tau\) and compute from the public key

\[
u=\operatorname{HighBits}(-c_0gb_1 2^d,m).
\]

For each candidate message \(M\), compute \(\mu=H_\mu(tr,M)\) and
\(\widetilde c=H_{ch}(\mu,u)\). If
\(\operatorname{SampleInBall}(\widetilde c)=c_0\), output the canonical
encoding of \((\widetilde c,z=0,h=0)\).

The verifier recomputes \(c=c_0\) and \(v=-c_0gb_1 2^d\). Since
\(\operatorname{UseHint}(0,v,m)=\operatorname{HighBits}(v,m)\), it
obtains \(u\) and reconstructs the supplied \(\widetilde c\). The
response and hint bounds hold. Therefore a successful trial is an
accepted signature obtained without any signing query.

In the ideal challenge-expansion model, a fresh trial succeeds with
probability \(1/|C_\tau|\). For a fixed deterministic expander, the
exact probability is the mass of \(c_0\) under a uniformly sampled
challenge seed; uniformity is the modeling assumption behind the stated
entropy. This construction also applies to the implementation's
alternative sign-stream map.

\subsection{4.2 Work bounds and units}\label{work-bounds-and-units}

Classical search needs approximately \(|C_\tau|\) trials. Grover search
reduces this to order \(\sqrt{|C_\tau|}\) predicate queries {[}G96{]}.
The entropy calculation in {[}E1{]} gives:

{\def\LTcaptype{none} % do not increment counter
\begin{longtable}[]{@{}
  >{\raggedright\arraybackslash}p{(\linewidth - 6\tabcolsep) * \real{0.2500}}
  >{\raggedleft\arraybackslash}p{(\linewidth - 6\tabcolsep) * \real{0.2500}}
  >{\raggedleft\arraybackslash}p{(\linewidth - 6\tabcolsep) * \real{0.2500}}
  >{\raggedleft\arraybackslash}p{(\linewidth - 6\tabcolsep) * \real{0.2500}}@{}}
\toprule\noalign{}
\begin{minipage}[b]{\linewidth}\raggedright
Profile
\end{minipage} & \begin{minipage}[b]{\linewidth}\raggedleft
\(\log_2\lvert C_\tau\rvert\)
\end{minipage} & \begin{minipage}[b]{\linewidth}\raggedleft
Classical trial exponent
\end{minipage} & \begin{minipage}[b]{\linewidth}\raggedleft
Grover query exponent
\end{minipage} \\
\midrule\noalign{}
\endhead
\bottomrule\noalign{}
\endlastfoot
128 & 131.579934 & 131.579934 & 65.789967 \\
256 & 257.007637 & 257.007637 & 128.503819 \\
512 & 512.004143 & 512.004143 & 256.002071 \\
\end{longtable}
}

The 128 profile is 14.21 bits below an 80-bit quantum-\emph{query}
target. A comparison with an 80-bit logical-\emph{gate} target requires
a reversible circuit, success-probability convention, memory model, and
depth accounting. None is supplied here. In particular, a query exponent
is not an unconditional gate-count break.

The 256 and 512 profiles exceed their corresponding nominal targets by
only 1.008 and 0.004 classical trial bits, or 0.504 and 0.002
Grover-query bits. These are single-route differences, not certified
margins against every attack. The longer stored challenge seed does not
change the search space used by this construction.

The theorem's \(2^{-|\widetilde c|}\) term can bound an unqueried
random-oracle guess. It does not bound the attack above, which makes the
relevant query and searches for a matching expanded challenge.

For illustration, \(\tau=21\) is the first weight giving at least 80
Grover-query bits at \(n=1024\); it gives 82.616220 bits. This is not a
validated replacement profile. At level 128 it also raises \(\beta\)
from 64 to 84, the two-test expected repetition count from 2.129771 to
2.692582, and the stated EUF residual bound from 393,217 to 434,177
{[}E1{]}. Compression and hint behavior must be analyzed again.

\section{5. A 344-bit descriptor for an equivalent SM3 signing
secret}\label{a-344-bit-descriptor-for-an-equivalent-sm3-signing-secret}

\subsection{5.1 Byte layout}\label{byte-layout}

The compact secret-XOF encoding in the SM3 implementations has the
following form before the outer counter:

{\def\LTcaptype{none} % do not increment counter
\begin{longtable}[]{@{}lr@{}}
\toprule\noalign{}
Field & Bytes \\
\midrule\noalign{}
\endhead
\bottomrule\noalign{}
\endlastfoot
Fixed prefix, role, level, rank, output length, seed length & 11 \\
Secret seed \(\rho'\) & 128 \\
Polynomial index & 1 \\
Lane index, zero for \(s_1\) & 1 \\
Total & 141 \\
\end{longtable}
}

The first 128 input bytes are two complete SM3 blocks. They contain the
fixed header and the first 117 bytes of \(\rho'\), and are shared by
every polynomial expansion in one profile. The polynomial index occurs
later.

Let \(V\) be the 256-bit chaining state after those two blocks. Every
remaining secret-dependent input byte is in the final 11 bytes of
\(\rho'\). Hence

\[
D=(V,\rho'[117:128])
\]

is a descriptor of length \(256+88=344\) bits from which the complete
\(s_1\) can be generated. Polynomial indices, lane identifiers, counter
values, total input lengths, and padding are fixed or public. This is an
exact data-flow property of the implemented construction.

\subsection{5.2 Equivalent-key search}\label{equivalent-key-search}

An adversary can enumerate all descriptors \(D\), continue SM3 from the
candidate state, and decode the resulting streams into a candidate
\(s_1\). With \(A\) regenerated from the public seed, it can compute the
candidate rounded \(b\), compress to \(b_1\), and compare with the
target public key.

The actual descriptor is included in this finite search. A candidate
yielding the public key supplies the same KeyGen relations: derive
\(s_2=As_1-gb\) and \(b_0=b-b_1 2^d\), retain the public seed and
\(tr\), and choose a new signing-randomness key \(K\). Verification does
not test the original \(K\) or require that the new secret fields arose
from the original KeyGen seed.

Thus an equivalent signing key has a search bound of at most \(2^{344}\)
descriptor evaluations, or order \(2^{172}\) Grover predicate queries.
Candidate evaluation includes hash expansion and public-key arithmetic.
The original key-generation seed, public-seed derivation, and stored
secret-key bytes need not be recovered.

This bound is incompatible with the 512-bit classical / 256-bit quantum
target under the corresponding search conventions. It is separate from,
and more expensive than, the SM3 message-collision forgery. It does not
apply this particular descriptor calculation to SHAKE.

The result is a search-space bound on an equivalent signing key. It is
not a measurement of the entropy of every stored secret-key field, and
no unknown-key recovery experiment is claimed. The relevant source
locations are the compact encoding in \emph{hash\_domain.h}, secret
expansion in \emph{arith/uniform.c}, and \emph{pseudoXOF} in
\emph{utils/auxfunc.c} {[}E9{]}. The executable check {[}E11{]}
reconstructs every secret-expansion byte from this descriptor for a
known test seed in all six SM3 implementation/profile combinations.

\section{6. Polynomial solutions of the relaxed 512-level SIS
instances}\label{polynomial-solutions-of-the-relaxed-512-level-sis-instances}

\subsection{6.1 The estimated problems}\label{the-estimated-problems}

The specification's Appendix A passes ordinary homogeneous infinity-norm
SIS instances to lattice-estimator. Their dimensions and single bounds
are

\[
(N,3N,q,B_{\mathrm{EUF}}),\qquad (N,2N,q,B_{\mathrm{SUF}}),
\]

where

\[
B_{\mathrm{EUF}}=\max\{\gamma_1-\beta,\;2\gamma_2+1+g\tau 2^{d-1}\},
\]

\[
B_{\mathrm{SUF}}=\max\{2(\gamma_1-\beta),\;4\gamma_2+2\}.
\]

At level 512,

\[
N=5120,\quad B_{\mathrm{EUF}}=4,947,969,\quad
B_{\mathrm{SUF}}=4,194,306.
\]

Both bounds exceed \(q/4=4,194,304\).

\subsection{6.2 A binary-linear-algebra
construction}\label{a-binary-linear-algebra-construction}

\textbf{Proposition.} Let \(q\) be divisible by four. For a matrix \(M\)
with at least \(2N\) columns, suppose its first two \(N\times N\) blocks
\(B,C\) are invertible modulo 2. A nonzero vector \(X\) with
\(MX=0\bmod q\) and \(\|X\|_\infty=q/4\) can be constructed in
polynomial time.

\textbf{Proof.} Put \(B_2=B\bmod2\) and \(C_2=C\bmod2\). Solve

\[
B_2u=C_2\mathbf 1
\]

over \(\mathbb F_2\) and lift \(u\) to a binary integer vector. Then
\(Bu+C\mathbf1\) is even. Define

\[
d_2=(Bu+C\mathbf1)/2\bmod2
\]

and solve \(C_2v=d_2\), again lifting \(v\) to a binary vector. The
vector

\[
w=(u,\mathbf1-2v,0,\ldots,0)
\]

satisfies

\[
Mw=Bu+C\mathbf1-2Cv=0\bmod4.
\]

Its second block consists of \(\pm1\), so it is nonzero and has infinity
norm one. Set \(X=(q/4)w\). Then \(MX=0\bmod q\) and
\(\|X\|_\infty=q/4\). This uses two linear solves over \(\mathbb F_2\)
and elementary integer operations. \(\square\)

For a uniformly random binary square matrix, the invertibility
probability is

\[
\prod_{j=1}^{N}(1-2^{-j})\longrightarrow 0.288788\ldots.
\]

The two-block event has limiting probability approximately 0.0834. It is
publicly testable. Consequently, the construction succeeds in polynomial
time on a constant fraction of uniformly random SIS instances. A
constant success probability is sufficient to invalidate a cryptographic
hardness claim for this distribution.

The result applies to both Appendix A instances at level 512. Extra
columns are assigned zero.

\subsection{6.3 Evidence and structured
extension}\label{evidence-and-structured-extension}

The artifact {[}E5{]} records an \(N=5120\) run with a norm-4,194,304
witness. It checks the relation modulo 4 and a random full-modulus lift
modulo \(2^{24}\). The recorded run took 10.56 seconds including
generation, solving, checking, and storage; this is one illustrative
run, not a general benchmark. It samples each block until the public
invertibility condition holds. The probability claim comes from the
formula above, not from this conditioned experiment. The saved witness
can be checked independently without trusting the reported timing.

For the structured relation \(A\Delta z-\Delta r=0\), assume
multiplication by \(A\bmod2\) is invertible. Let \(u\) be the binary
lift of \(A_2^{-1}\mathbf1\), and let \(r\) be the centered lift of
\(Au\bmod4\). Every coefficient of \(r\) is odd, hence in \(\{-1,1\}\).
Then

\[
(\Delta z,\Delta r)=(q/4)(u,r)
\]

solves the relation with both block norms at most \(q/4\). For a
uniformly sampled \(A\) at \(k=5\), the unit probability is \(15/32\),
as derived in Section 8.1.

\subsection{6.4 Why this does not yet forge a
signature}\label{why-this-does-not-yet-forge-a-signature}

The actual SUF extraction on pp.~31--33 retains separate conditions:

\[
\|\Delta z\|_\infty<2(\gamma_1-\beta),\qquad
\|\Delta r\|_\infty\leq4\gamma_2+2.
\]

At level 512 the strict response threshold is 4,193,956. It is 348 below
\(q/4\); the largest permitted integer coefficient is 4,193,955, or 349
below \(q/4\). The structured witness has a response coefficient equal
to \(q/4\) and therefore fails this condition. Even a short solution of
the extracted relation would still need to be turned into accepting
signatures with consistent commitments, hashes, and hints.

The EUF relation additionally restricts \(c\) to an expanded sparse
challenge and binds the target to the random oracle. Assigning zero to a
relaxed challenge block does not satisfy that relation.

The established conclusion is that the \emph{relaxed SIS instances used
in the estimates are easy}. Hardness of the more constrained signature
relations remains unestablished. The 691/692-bit figures in Table 2
cannot be interpreted as hardness of the relaxed problems, even if they
correctly price one lattice-reduction attack.

All three profiles lie near analogous powers-of-two boundaries:

{\def\LTcaptype{none} % do not increment counter
\begin{longtable}[]{@{}
  >{\raggedright\arraybackslash}p{(\linewidth - 6\tabcolsep) * \real{0.2500}}
  >{\raggedleft\arraybackslash}p{(\linewidth - 6\tabcolsep) * \real{0.2500}}
  >{\raggedleft\arraybackslash}p{(\linewidth - 6\tabcolsep) * \real{0.2500}}
  >{\raggedleft\arraybackslash}p{(\linewidth - 6\tabcolsep) * \real{0.2500}}@{}}
\toprule\noalign{}
\begin{minipage}[b]{\linewidth}\raggedright
Profile
\end{minipage} & \begin{minipage}[b]{\linewidth}\raggedleft
Scale
\end{minipage} & \begin{minipage}[b]{\linewidth}\raggedleft
Strict response-difference threshold
\end{minipage} & \begin{minipage}[b]{\linewidth}\raggedleft
Residual bound
\end{minipage} \\
\midrule\noalign{}
\endhead
\bottomrule\noalign{}
\endlastfoot
128 & \(q/32=524,288\) & \(<524,160\) & 524,290 \\
256 & \(q/16=1,048,576\) & \(<1,048,432\) & 1,048,578 \\
512 & \(q/4=4,194,304\) & \(<4,193,956\) & 4,194,306 \\
\end{longtable}
}

The modulo-4 construction does not automatically extend to the modulo-16
or modulo-32 problems needed for the first two rows. Those rows identify
boundaries requiring analysis, not attacks on the 128 or 256 profiles.

\section{7. What the security proof still
needs}\label{what-the-security-proof-still-needs}

\subsection{7.1 The assumed relations are more specific than generic
SIS}\label{the-assumed-relations-are-more-specific-than-generic-sis}

Theorem 4, pp.~31--32, assumes hardness of a self-targeted relation.
Given \((A,b)\), the adversary must produce
\((\mu,z,\widetilde c,c,r,w_1)\) satisfying

\[
Az-cgb-r=mw_1,\qquad
c=\operatorname{SampleInBall}(\widetilde c),\qquad
\widetilde c=H_{ch}(\mu,w_1),
\]

together with the response, challenge, and residual constraints. This is
neither ordinary homogeneous SIS nor a target problem with an
independently uniform right-hand side. Theorem 5 also retains two
separate norm bounds in the homogeneous relation
\(A\Delta z-\Delta r=0\).

Appendix A estimates homogeneous SIS with \(N\) rows, \(3N\) or \(2N\)
columns, and one bound. Pricing lattice reduction on that relaxation
does not establish hardness of the actual assumption. In particular, an
efficient solution of the relaxation need not yield a forgery, while
failure of one attack on the relaxation does not imply that all attacks
on the original relation are expensive. Section 6 supplies an explicit
omitted attack on the relaxation itself.

The key-generation hybrid uses \emph{decision} RR-LWR with the submitted
small-secret distribution. A search-LWE estimate does not establish that
decisional assumption without a reduction applicable to the ring, basis,
distribution, sample count, and rounding parameters.

\subsection{7.2 An explicit signing simulator is possible in
principle}\label{an-explicit-signing-simulator-is-possible-in-principle}

The proof refers to a public transcript distribution without specifying
an executable sampler or its acceptance probability. That omission
should be repaired constructively.

The reduction receives the full \((A,b)\) and can compute \(b_0\) even
though it publishes only \(b_1\). The following is a candidate simulator
for a fresh idealized signing attempt with a fixed message
representative \(\mu\):

\begin{enumerate}
\def\labelenumi{\arabic{enumi}.}
\tightlist
\item
  Sample \(\widetilde c\) uniformly and expand \(c\).
\item
  Sample \(z\) uniformly in the strict response box
  \(\|z\|_\infty<\gamma_1-\beta\).
\item
  Set \(t=Az-cgb\) and \(\delta=cgb_0\). Reject unless
  \(\|\operatorname{LowBits}(t,m)\|_\infty<\gamma_2-\beta\) and
  \(\|\delta\|_\infty<\gamma_2\).
\item
  Set \(v=t+\delta\) and \(h=\operatorname{MakeHint}(-\delta,v,m)\).
  Reject if \(\operatorname{HW}(h)>\omega\).
\item
  Set \(w_1=\operatorname{HighBits}(t,m)\) and program a fresh
  \(H_{ch}(\mu,w_1)\) to \(\widetilde c\).
\end{enumerate}

For a real key, write \(Y=z-cs_1\). The response box and
\(\|cs_1\|_\infty\leq\beta\) place \(Y\) in the original mask box. Also,

\[
AY+cs_2=A(z-cs_1)+c(As_1-gb)=t.
\]

Thus, for one fresh idealized attempt, the change of variables
\((Y,c)\leftrightarrow(z,c)\) matches the tests and hints, after
conditioning on the response test. This observation does not require
\(AY\) for bounded \(Y\) to be uniform modulo \(q\). It gives a
plausible route to a proof; it is not a completed adaptive simulation.

A full argument must handle rejected attempts, prior oracle queries,
repeated messages, the seeded mask stream and its lane indices, counter
exhaustion, and the acceptance probability after replacing \(b\) by
uniform. One useful proof order is to establish the simulator in the
real-key distribution and then analyze the RR-LWR hybrid. The
specification must make clear which distributions and oracle interfaces
each step uses.

\subsection{7.3 Uninstantiated statistical
losses}\label{uninstantiated-statistical-losses}

Theorem 4 introduces \(\Delta_{\mathrm{sim}}\) and a commitment
point-probability bound \(2^{-\alpha}\), but supplies no concrete
values. The resulting loss includes

\[
Q_S\Delta_{\mathrm{sim}}+Q_S(Q_H+Q_S)2^{-\alpha}.
\]

These terms cannot be omitted when assigning concrete security levels.
For example, if \(Q_S=Q_H=2^{64}\), making the programming term alone at
most \(2^{-\lambda}\) requires

\[
\alpha\geq\lambda+129.
\]

The required values are 257, 385, and 641 for the three nominal
classical targets. Making the simulation term alone at most
\(2^{-\lambda}\) requires
\(\Delta_{\mathrm{sim}}\leq2^{-(\lambda+64)}\). These are illustrative
necessary budgets for the individual terms, not a complete security
calculation.

A basic counting argument can help with a commitment bound. For fixed
\(A\), let \(K_A=|\ker(A:S_q\to S_q)|\). If \(Y\) is uniform in the full
mask box, then each value of \(AY\) has at most \(K_A\) preimages in
that box. A high-bit vector has \(m^N\) possible low-bit lifts. Since
\(\gamma_1=m\),

\[
\max_u\Pr[\operatorname{HighBits}(AY,m)=u]
\leq\frac{K_A m^N}{(2\gamma_1)^N}
=K_A2^{-N}.
\]

Conditioning on an acceptance event of probability \(a\) can enlarge
this bound by a factor \(1/a\). This is a starting point for analysis,
not a bound on every adaptive simulator transcript. It requires control
of the kernel, acceptance probability, and conditioning on oracle
history.

The theorem defines worst-case quantities over keys in the post-RR-LWR
hybrid. Degenerate keys, such as \(A=b=0\), show why an unrestricted
uniform high-entropy claim cannot hold: the candidate commitment is then
fixed. Such keys may have negligible probability, but a proof must
explicitly isolate bad-key events and charge their probabilities.

\subsection{7.4 The public seed remains
observable}\label{the-public-seed-remains-observable}

The public key contains \(\rho\), from which an adversary can recompute
\(A\). Replacing the expanded \(A\) by an independently uniform element
while retaining an unprogrammed \(\rho\) is distinguishable by that
recomputation.

A seeded-sampler hybrid needs an explicit interface: for example,
programming the selected public-expansion oracle calls and bounding the
probability that those inputs were queried earlier. Merely assigning a
symbol \(\Delta_A\) to this step does not quantify its cost. Exact
uniformity of coefficients produced by an ideal stream, which is
particularly simple for a power-of-two modulus, is a different issue
from consistency with the public seed.

\subsection{7.5 LWR-to-LWE comparisons require their actual
hypotheses}\label{lwr-to-lwe-comparisons-require-their-actual-hypotheses}

The bounded-error LWR result of Bogdanov et al.~{[}BGM16, Theorems
1--2{]} uses an auxiliary error bound \(B\) and the strict condition
\(q>2pB\). Its comparison includes a factor of the form

\[
(1+2pB/q)^{m_s},
\]

as well as a square of the relevant success probability. Here \(m_s\)
counts scalar coordinates in the comparison. The parameter \(B\) is not
automatically the scheme's \(\eta\) or \(\beta\).

For the submitted parameters, setting \(B=\eta=q/(2p)\) violates the
strict inequality. Setting \(B=1\) satisfies it, but changes the error
distribution being compared. For illustration, taking \(m_s=N\) gives

\[
\log_2(1+2p/q)^N=N\log_2(1+1/\eta),
\]

which is approximately 329.65, 1198.00, and 2995.01 bits. These figures
describe the size of this factor under that substitution; they are not
standalone security losses for a fully instantiated radical-ring
reduction.

The convenient sufficient condition \(q\geq2m_sBp\) keeps this factor at
most \(e\); failure of that sufficient condition does not invalidate
every possible reduction. The ring theorem additionally assumes a secret
supported on units. The submitted secret distribution has substantial
nonunit mass, and its coefficient representation must be matched to the
theorem. The required decisional, distributional, and basis arguments
remain to be supplied.

\subsection{7.6 Classical ROM, deterministic signing, and stronger
properties}\label{classical-rom-deterministic-signing-and-stronger-properties}

The specification explicitly limits Theorems 4--5 to the classical
random-oracle model. Its quantum attack estimates are not a QROM proof.
A revised document should distinguish a conjectured post-quantum work
factor from a proven QROM reduction.

The signing interface also permits deterministic operation. The
specification correctly states that the randomized signing oracle used
in its proof excludes this case. Repeated deterministic signing of the
same message must repeat the signature, with the corresponding oracle
and seed correlations. Claims covering that interface require a separate
argument.

Theorem 6's collision-based reasoning for stronger binding properties is
useful in the stated ideal model. Its instantiation still needs the
actual primitive collision bounds and any high-min-entropy assumptions
used for the weak non-resignability experiment. Neither a 1024-bit
digest length nor the existence of a classical ROM unforgeability
argument discharges those assumptions. No two-key exclusive-ownership
attack is established in this document.

\section{8. Ring structure, basis geometry, and challenge
algebra}\label{ring-structure-basis-geometry-and-challenge-algebra}

\subsection{8.1 Reduction modulo 2 and
nonunits}\label{reduction-modulo-2-and-nonunits}

Eliminating \(y\) gives

\[
S_q\cong(\mathbb Z/q\mathbb Z)[X]/\big((X^k-2)^{1024}+1\big).
\]

Modulo 2 the defining polynomial factors as follows:

\[
\begin{aligned}
k=1:\quad &(X+1)^{1024},\\
k=2:\quad &(X+1)^{2048},\\
k=5:\quad &(X+1)^{1024}
  (X^4+X^3+X^2+X+1)^{1024}.
\end{aligned}
\]

The quartic is irreducible over \(\mathbb F_2\). An element modulo
\(2^{24}\) is a unit exactly when its reduction modulo 2 is coprime to
these distinct factors. Hence the unit probabilities for a uniformly
random element are

\[
\frac12,\quad\frac12,\quad
(1-2^{-1})(1-2^{-4})=\frac{15}{32}.
\]

The same probabilities hold in the ideal independent-coefficient model
for \(s_1\), because its coefficient intervals give unbiased independent
parity bits. These facts are verified in {[}E2, E6{]}.

For a fixed nonunit \(s_1\), a uniform \(A\) gives \(As_1\) in the
proper ideal generated by \(s_1\), rather than uniformly over the ring.
Thus an unconditional uniform-product heuristic is false. This does not
itself refute RR-LWR: rounding can obscure ideal membership, and the
assumption concerns rounded samples.

\subsection{\texorpdfstring{8.2 A decomposition of the \(k=5\)
ring}{8.2 A decomposition of the k=5 ring}}\label{a-decomposition-of-the-k5-ring}

In \(R_q\), set \(t=y^{1229}\). Since \(5\cdot1229=1\bmod2048\) and
\(y^{2048}=1\), one has \(t^5=y\). Therefore

\[
a=x-t,\qquad
b=x^4+tx^3+t^2x^2+t^3x+t^4
\]

satisfy the exact identity \(ab=x^5-t^5=2\). In particular,
\(a^{24}b^{24}=0\) modulo \(2^{24}\).

There is also a coprime factorization. Hensel lifting starts from
\(t\bmod2\) and produces \(r\in R_q\) such that \(r^5=y+2\), since the
derivative \(5t^4\) is a unit. Then

\[
x^5-y-2=(x-r)
(x^4+rx^3+r^2x^2+r^3x+r^4).
\]

The quartic evaluated at \(x=r\) is \(5r^4\), a unit, so the factors are
comaximal. The Chinese remainder theorem decomposes \(S_q\) into
components of ranks 1024 and 4096 over \(\mathbb Z/q\mathbb Z\).

This is an exact algebraic decomposition. Its usefulness for an attack
depends on what it does to the small-secret distribution and coefficient
norm. In the diagnostic {[}E6{]}, projecting one vector with original
coefficients in \([-2,1]\) produced a maximum centered coefficient of
8,386,517, close to \(q/2=8,388,608\). The explicit annihilators above
had norms 8,370,000 and 8,260,608. Neither example supplies a short
response satisfying the signature bounds.

These observations rule out using those particular projected vectors as
immediate attacks. They do not rule out other projections, smaller
annihilators, multilevel lifting, or attacks exploiting the two
components jointly.

\subsection{8.3 Even-weight challenges are
nonunits}\label{even-weight-challenges-are-nonunits}

Over \(\mathbb F_2\), signs disappear and a weight-\(\tau\) challenge
satisfies

\[
c(1)=\tau\bmod2.
\]

At levels 128 and 256, every challenge has even weight, so it is
divisible by \(y+1\) and is a nonunit. Multiplication by such a
challenge has a nontrivial kernel. At level 512, the odd weight makes
\(c\) a unit in \(R_q\) and hence in \(S_q\).

Nonunit challenges are not, by themselves, a forgery. The extraction
displayed in the specification does not simply invert \(c\).
Nevertheless, this deterministic restriction matters for distribution
claims and for any argument involving multiplication by a challenge or
challenge difference. Changing the parity of \(\tau\) would require a
fresh parameter and rejection analysis; it would not repair the other
issues.

\subsection{8.4 Maximality should be proved for the submitted
parameters}\label{maximality-should-be-proved-for-the-submitted-parameters}

For \(f(X)=(X^k-2)^n+1\), direct resultant calculation gives

\[
|\operatorname{disc}(f)|=(nk)^{nk}(2^n+1)^{k-1}.
\]

The defining polynomial is irreducible for the relevant radical-ring
family {[}BRVW26, Lemma 1{]}. To identify \(\mathbb Z[X]/(f)\) with the
full ring of integers, one must additionally exclude primes dividing the
index.

For \(n=1024\), \(k\in\{1,2,5\}\), the calculations in {[}E2{]}
establish the required local checks at 2 and, for \(k=5\), at 5. At a
prime dividing \(2^{1024}+1\), with \(p\nmid nk\), the Dedekind
criterion reduces to whether \(p^2\) divides that Fermat number. Its
published complete factorization contains four distinct prime factors
{[}B99{]}. Together, these facts establish maximality for the submitted
profiles. The code checks the arithmetic of the factors; the primality
of the largest factor relies on the published factorization.

A general claim for arbitrary ranks needs additional hypotheses. For
example, take \(n=2\), \(k=17\), and

\[
f=(X^{17}-2)^2+1,\qquad g=(X-2)^2+1.
\]

Modulo 17, \(f=g^{17}\). For \(T=(f-g^{17})/17\),

\[
T\bmod(17,g)=8+10X,\qquad
\gcd(T,g)\bmod17=X-6.
\]

Since \(g\) is squarefree modulo 17, Dedekind's index criterion implies
that 17 divides the index. The exact calculation and irreducibility
check appear in {[}E10{]}. This is a counterexample to unrestricted
extrapolation; it is not an index defect in any submitted profile.

Even when the global index question is settled, maximality alone is not
a hardness theorem. For arithmetic modulo \(2^{24}\), the local analysis
at 2 is the directly relevant order check.

\subsection{8.5 Exact embedding
distortion}\label{exact-embedding-distortion}

Let \(\zeta_j=\exp((2j+1)\pi i/n)\). In the coefficient basis
\(y^a x^b\), with \(0\leq a<n\) and \(0\leq b<k\), the singular values
of the full complex evaluation matrix are

\[
\sqrt{nk}\,|2+\zeta_j|^{b/k},
\qquad 0\leq j<n,\quad0\leq b<k.
\]

To see this, group embeddings by \(\zeta_j\) and by the \(k\) choices of
the radical root. Fourier orthogonality in the root index separates the
\(b\) blocks; Fourier diagonalization in the \(y\) index then gives the
displayed values. This uses the usual Euclidean norm on all complex
embeddings, equivalent to the standard \(\sqrt2\)-weighted real
representation.

The condition number is

\[
\kappa=(5+4\cos(\pi/n))^{(k-1)/(2k)}.
\]

At \(n=1024\) it is approximately 1, 1.732049902, and 2.408222670. Small
numerical matrix checks are recorded in {[}E10{]}. The distortion is
modest, but the transport is not isometric for \(k>1\); both covariance
and volume factors must be accounted for in a reduction using canonical
geometry. A generic estimator call does not incorporate them merely by
setting the dimension to \(nk\).

For \(k=5\), the field is also not Galois over \(\mathbb Q\). Normality
would require all five radical roots and hence \(\zeta_5\), alongside
the existing \(\zeta_{2048}\). The compositum has degree 4096, which
cannot divide the field degree 5120. Cyclotomic Ring-LWE arguments using
a full automorphism group therefore cannot be imported without checking
their hypotheses. Complex conjugation still exists; non-Galois does not
mean the absence of every symmetry.

\section{9. Public transcript information and functional
checks}\label{public-transcript-information-and-functional-checks}

\subsection{9.1 Signatures give bounded observations of the omitted key
part}\label{signatures-give-bounded-observations-of-the-omitted-key-part}

An observer can reconstruct \(c\) and \(w_1\) from a valid signature and
calculate

\[
e=\operatorname{ctr}_q(Az-cgb_1 2^d-mw_1).
\]

For an honest accepted signature, put
\(r_0=\operatorname{LowBits}(Az-cgb,m)\). The signing tests give

\[
e=cgb_0+r_0,\qquad
\|r_0\|_\infty<\gamma_2-\beta,\qquad
\|cgb_0\|_\infty<\gamma_2.
\]

There is no centered-modulus ambiguity because \(2\gamma_2-\beta<q/2\).
Thus multiple signatures expose bounded linear observations of \(b_0\).

The noise is affected by rejection and by the other transcript
variables. It must not be treated as independent Gaussian noise without
justification. A useful next experiment would measure the conditional
distributions and attempt recovery of \(b_0\) in controlled, known-key
tests.

Recovering \(b_0\) would reveal the full rounded value \(b\), which is
already public in the stated RR-LWR assumption. It would not
automatically recover \(s_1\) or supply a forgery. The purpose of this
observation is to identify a concrete transcript-analysis problem and
avoid silently treating the omitted key part as permanently hidden.

\subsection{9.2 Results that do check
out}\label{results-that-do-check-out}

The following arithmetic is internally consistent with the parameter
macros and encodings {[}E1{]}:

{\def\LTcaptype{none} % do not increment counter
\begin{longtable}[]{@{}lrrr@{}}
\toprule\noalign{}
Quantity & 128 & 256 & 512 \\
\midrule\noalign{}
\endhead
\bottomrule\noalign{}
\endlastfoot
Public-key bytes & 1344 & 2368 & 5184 \\
Secret-key bytes & 2432 & 4608 & 11776 \\
Signature bytes & 2564 & 5449 & 14713 \\
Encoded hint bytes & 100 & 265 & 505 \\
Two-test expected repetitions & 2.129771 & 2.338962 & 3.590045 \\
\(B_{\mathrm{EUF}}\) & 393,217 & 1,114,113 & 4,947,969 \\
\(B_{\mathrm{SUF}}\) & 524,290 & 1,048,578 & 4,194,306 \\
\end{longtable}
}

The repetition calculation covers the stated response and low-bit model.
It is not a measurement of the complete signer with all rejection tests
and correlations.

The key-generation residual bound and the accepted-signature
verification identities are consistent with the centered rounding
conventions. Exhaustive small-modulus checks in {[}E2{]} verify hint
recovery, high-bit stability under the stated strict inequalities, and
the residual bound for both hint values. The SUF argument's
nonzero-extraction step is consistent with canonical encodings and
uniqueness of the hint when the high-bit moduli are 64, 32, and 8.

The two NTT primes have product

\[
1,152,837,945,367,908,353.
\]

Against the supplied worst-case CRT reconstruction bounds, the
corresponding headroom factors are approximately 255.98, 32.00, and 2.46
{[}E1{]}. This verifies the integer reconstruction range used in that
calculation. It is not a proof that every implementation intermediate
avoids overflow.

These positive checks matter: the unsupported security claims should not
be conflated with a demonstrated failure of the basic signing identity
or an incorrect byte count.

\section{10. Specification, implementation, and estimate
consistency}\label{specification-implementation-and-estimate-consistency}

\subsection{10.1 A reproducible disagreement in the challenge
map}\label{a-reproducible-disagreement-in-the-challenge-map}

Algorithm 2 on p.~22 requests \(\lceil\tau/8\rceil\) bytes for the sign
stream: 2, 5, and 11 bytes. All nine implementation trees request 11
bytes. The requested output length is part of the hashed domain
encoding. Consequently, requesting 11 bytes and then taking a prefix
changes the hash input; it is not equivalent to requesting the shorter
output.

The source check {[}E7{]} covers the three levels in each of the
reference, optimized, and additional families. The C regression {[}E8{]}
executes the submitted wrappers in all nine combinations with an
all-zero public challenge seed. The SM3 results, identical in the
reference and optimized families, are:

{\def\LTcaptype{none} % do not increment counter
\begin{longtable}[]{@{}
  >{\raggedright\arraybackslash}p{(\linewidth - 6\tabcolsep) * \real{0.2500}}
  >{\raggedleft\arraybackslash}p{(\linewidth - 6\tabcolsep) * \real{0.2500}}
  >{\raggedleft\arraybackslash}p{(\linewidth - 6\tabcolsep) * \real{0.2500}}
  >{\raggedleft\arraybackslash}p{(\linewidth - 6\tabcolsep) * \real{0.2500}}@{}}
\toprule\noalign{}
\begin{minipage}[b]{\linewidth}\raggedright
Profile
\end{minipage} & \begin{minipage}[b]{\linewidth}\raggedleft
Specified bytes
\end{minipage} & \begin{minipage}[b]{\linewidth}\raggedleft
Implementation bytes
\end{minipage} & \begin{minipage}[b]{\linewidth}\raggedleft
Differing used sign bits
\end{minipage} \\
\midrule\noalign{}
\endhead
\bottomrule\noalign{}
\endlastfoot
128 & 2 & 11 & 5 \\
256 & 5 & 11 & 18 \\
512 & 11 & 11 & 0 \\
\end{longtable}
}

At level 128 the specification-length stream is hexadecimal \emph{3b80},
while the corresponding prefix of the implementation stream is
\emph{9bc5}. At level 256 the two values are \emph{120021c7c5} and
\emph{82543e7c36}. The index stream is unchanged, so these sign
differences alter the expanded challenge itself.

The additional SHAKE family exhibits the same disagreement: \emph{6da1}
versus \emph{3930} at level 128, and \emph{40f92b5756} versus
\emph{3598e19826} at level 256. There are respectively 6 and 18
differing used sign bits. At level 512, the specified and implemented
lengths coincide in all families.

This disagreement does not establish lower entropy for either map. It
establishes that the PDF and implementations define different schemes at
two levels. Agreement among implementations or known-answer tests
sharing the same map would not resolve that discrepancy.

\subsection{10.2 Other normative ambiguities and interface
issues}\label{other-normative-ambiguities-and-interface-issues}

Several smaller differences should be resolved in a normative revision:

\begin{itemize}
\tightlist
\item
  \emph{UseHint} treats a zero low part differently: the PDF's branch
  for \(r_0\leq0\) decrements, while the code increments when \(r_0=0\).
  Honest nonzero hints avoid this boundary under the signing bounds. An
  arbitrary candidate signature can contain such a hint, so the
  verification language still needs one definition.
\item
  The document defines a generic \emph{Pack} operation for keys and
  signatures. It should explicitly require \(H_{ch}\) to receive
  \(\operatorname{Pack}(w_1,\log_2(q/m))\), or state another exact
  encoding. The coefficient widths are 6, 5, and 3 bits. The current
  hash notation leaves this connection to be inferred.
\item
  The deterministic signing option needs a matching algorithm and
  security statement, as discussed in Section 7.6.
\item
  Challenge-refill limits need one definition. Both the specification
  and wrapper encode the index-stream counter in 32 bits, but the code
  aborts after counter 255, having exhausted 256 buffers; Algorithm 2
  gives no such cap. Counter exhaustion or wraparound in sampling and
  signing needs specified behavior and a quantified failure probability.
\item
  The current domain header identifies the level and rank. It does not
  bind every parameter. This separates the submitted profiles, but
  future parameter changes under the same identifiers would require
  explicit version separation.
\end{itemize}

These items call for executable specification tests, including
adversarial encodings and boundary inputs, as well as honest-signature
test vectors.

\subsection{10.3 Randomness and side-channel
claims}\label{randomness-and-side-channel-claims}

The supplied SM3 random-generation wrapper uses global generator
selection and error-return interfaces. Some callers do not propagate
failure to the signing or key-generation API. A production interface
must make initialization, entropy acquisition, generator selection, and
error handling explicit.

A 55-byte seed length is not evidence of 440-bit or 512-bit security.
NIST SP 800-90A's 440-bit seed length for SHA-256 Hash\_DRBG is an
internal-state parameter; the standard separately specifies security
strength and approved hash instantiations {[}SP90A{]}. Substituting SM3
and reusing the seed length does not inherit an SP 800-90A validation
claim. This is an integration and assurance issue; no additional
mathematical forgery is inferred solely from the seed length.

A power-of-two modulus can simplify arithmetic and permit constant-time
implementations. It does not establish resistance to timing, cache,
power, electromagnetic, or fault attacks. Such claims need
implementation-specific evidence, particularly for secret expansion,
rejection loops, arithmetic, and error paths.

\subsection{10.4 Estimator output is not a security lower
bound}\label{estimator-output-is-not-a-security-lower-bound}

Appendix A uses \emph{lattice-estimator}, including a pinned source
revision {[}LE{]}. That revision's SIS implementation prices
lattice-reduction strategies; it does not test the modulo-4 construction
in Section 6. A large returned exponent is therefore compatible with an
easy instance through another algorithm.

The different numbers labeled Core-SVP, arithmetic operations, and
refined attack costs need not be numerically equal. A Core-SVP exponent
based on a block size omits algorithm-specific preprocessing and other
overhead. Such differences should be explained using the exact attack
variant, cost model, success probability, and operation convention,
rather than treated automatically as contradictions or combined
selectively.

The supplied cost-fitting range ends at block size 1024, while reported
512-level choices reach roughly 2181--2788. This is extrapolation. The
submission should publish the fitting data, residuals, extrapolation
rule, uncertainty, estimator revision, full inputs, and unedited
outputs. No fit can repair the use of an inapplicable problem model.

\subsection{10.5 Parameter and presentation
corrections}\label{parameter-and-presentation-corrections}

The following points also affect reproducibility or interpretation:

\begin{itemize}
\tightlist
\item
  Increasing \(d\) changes the public-key compression error and the
  stated EUF bound. For example, changing the 128-level value from 11 to
  12 changes \(B_{\mathrm{EUF}}\) from 393,217 to 524,289. A claim that
  this change has no security effect is too broad. Increasing \(\omega\)
  similarly enlarges the accepted hint set; its effect needs analysis.
\item
  \emph{Power2Round} directly compresses the public key. Any
  signature-size benefit is indirect and must be explained through the
  chosen parameters.
\item
  The appendix's generic signature-size expression using a fixed 64-byte
  challenge and omitting hints is not valid for all three profiles. The
  actual encoded sizes in Section 9.2 are consistent with the tables and
  code.
\item
  An expected number of \emph{additional} 64-word refill blocks must
  exclude the initial block and account for the discrete threshold.
  Under the stated sampling approximation the expectations are near 0,
  0, and 1, rather than simply the expected word count divided by 64.
\item
  Table 7 gives 204,968 signing cycles for Octarine-128 and 203,781 for
  its ML-DSA comparison, contrary to a blanket signing-speed advantage.
  The difference is about 0.58\%; differing compiler flags further limit
  the comparison. Performance claims need comparable builds, platform
  details, and dispersion across runs.
\end{itemize}

\section{11. Disposition and requirements for a
revision}\label{disposition-and-requirements-for-a-revision}

The submitted document should not be used to endorse the advertised
security levels. The 512-bit classical claim must be withdrawn for both
supplied hash families. The SM3 256-level implementation also has a
direct generic signature-transfer bound around \(2^{128}\).

{\def\LTcaptype{none} % do not increment counter
\begin{longtable}[]{@{}
  >{\raggedright\arraybackslash}p{(\linewidth - 2\tabcolsep) * \real{0.5000}}
  >{\raggedright\arraybackslash}p{(\linewidth - 2\tabcolsep) * \real{0.5000}}@{}}
\toprule\noalign{}
\begin{minipage}[b]{\linewidth}\raggedright
Family and profile
\end{minipage} & \begin{minipage}[b]{\linewidth}\raggedright
Principal established obstruction
\end{minipage} \\
\midrule\noalign{}
\endhead
\bottomrule\noalign{}
\endlastfoot
SM3, 128 & Challenge-search bound of \(2^{65.79}\) quantum queries;
unresolved proof and conformance issues \\
SM3, 256 & Message-collision signature transfer near \(2^{128}\)
classical work \\
SM3, 512 & Same collision bound; 344-bit equivalent-key descriptor; easy
relaxed SIS instances \\
SHAKE, 128 & Same challenge-search query bound and unresolved
proof/conformance issues \\
SHAKE, 256 & No below-target classical forgery established here; proof
and model obligations remain \\
SHAKE, 512 & Message-collision strength capped at 256 bits; easy relaxed
SIS instances \\
\end{longtable}
}

The quantum-query entry for level 128 must retain the gate-cost
qualification in Section 4.2. A profile without a below-target attack in
this document is not thereby validated.

A technically reviewable revision should satisfy the following
requirements:

\begin{enumerate}
\def\labelenumi{\arabic{enumi}.}
\tightlist
\item
  \textbf{Define the claim and the primitives.} State classical and
  quantum resource models, allowable query counts, target advantages,
  and per-role hash requirements. Specify each XOF and generator
  completely.
\item
  \textbf{Analyze the actual forgery relations.} Preserve the separate
  response and residual bounds, challenge support, target dependence,
  and oracle consistency. Include powers-of-two methods, modular
  projections, and structured kernels alongside lattice reduction.
\item
  \textbf{Complete the proof instantiation.} Give an executable
  simulator, acceptance bounds, commitment entropy, seeded-oracle
  programming, bad-key terms, and the applicable LWR reduction or an
  explicitly stated independent decisional assumption. Separate
  randomized ROM, deterministic, and QROM claims.
\item
  \textbf{Reconcile specification and code.} Provide one challenge map,
  one rounding and hint convention, exact encodings, versioned domains,
  failure semantics, and conformance tests.
\item
  \textbf{Re-estimate and reproduce.} Publish parameter derivations,
  complete estimator inputs and outputs, benchmark settings, and test
  artifacts. Reassess rejection rates and signature sizes after changes.
\item
  \textbf{Invite analysis of the revised construction.} Priority
  questions are constrained powers-of-two attacks near the response
  thresholds, multitranscript information about \(b_0\), componentwise
  attacks at \(k=5\), and the real distribution of accepted commitments.
\end{enumerate}

Further research on a corrected SHAKE-based construction at lower
claimed levels may be reasonable. It should proceed as an unvalidated
candidate with revised claims and a frozen specification. The available
evidence does not justify deployment or a claim that the current
submission has survived cryptanalysis at its advertised levels.

\section{12. Reproducibility and
evidence}\label{reproducibility-and-evidence}

\subsection{12.1 Evidence inventory}\label{evidence-inventory}

Paths below are relative to the distribution root. They identify
programs, source code, and data; none is a prerequisite for following
the mathematical constructions in Sections 3--8.

\textbf{{[}E1{]} Parameter arithmetic.}
\href{code/verify_parameters.py}{\emph{verify\_parameters.py}}
recomputes challenge entropies, byte counts, stated bounds, rejection
approximations, and CRT ranges using Python's standard library. Its
\href{evidence/verify_parameters.json}{recorded output} is included.

\textbf{{[}E2{]} Algebra and rounding checks.}
\href{code/verify_algebra.py}{\emph{verify\_algebra.py}} checks the
discriminant formula on small degrees, the relevant full-degree local
index calculations, Fermat-factor arithmetic, parity identities, and
exhaustive small-modulus hint properties. It does not independently
prove the primality of the largest published Fermat factor. See
\href{evidence/verify_algebra.json}{output}.

\textbf{{[}E3{]} Collision-propagation model.}
\href{code/verify_kdf_collision.py}{\emph{verify\_kdf\_collision.py}}
gives a toy state-collision experiment with a 12-bit state. Its reduced
state size is intentional; it is not a collision in full SM3. See
\href{evidence/verify_kdf_collision.txt}{output}.

\textbf{{[}E4{]} SM3 continuation.}
\href{code/verify_sm3_mechanism.c}{\emph{verify\_sm3\_mechanism.c}}
checks continuation against the supplied SM3 code. Any forced
internal-state equality is a mechanism test, not a found collision.
Results are in
\href{evidence/original_code.json}{\emph{original\_code.json}}.

\textbf{{[}E4a{]} Supplied-state 20-step collision.}
\href{code/verify_sm3_reduced_collision.c}{\emph{verify\_sm3\_reduced\_collision.c}}
replays the second-block pair from {[}MNS13, Table 3{]} with an
adaptation of the submitted compression function. It counts the
differing bytes, checks the published output and 64 common counter
continuations, and requires the unmodified 64-step function to give
different states. Four inputs cross-check the adapter against that
function. \href{evidence/verify_sm3_reduced_collision.txt}{Recorded
output} is also captured by the original-source driver in
\href{evidence/original_code.json}{\emph{original\_code.json}}. The
supplied chaining input and reduced step count preclude treating this as
an Octarine signature-transfer instance.

\textbf{{[}E4b{]} Reduced-state signature-transfer model.}
\href{code/signature_transfer/run_signature_transfer.py}{\emph{run\_signature\_transfer.py}}
generates a 48-bit state variant from the hash-verified submitted
utility, runs one chosen-message signing query, and checks transfer to a
never-signed equal-length message. The
\href{evidence/signature_transfer/campaign_48.json}{campaign
transcript},
\href{evidence/signature_transfer/probe_fullwidth.json}{full-width
controls},
\href{evidence/signature_transfer/public_artifact_verification.json}{independent
public-artifact checks}, and
\href{evidence/signature_transfer/scale.json}{finite-width scaling
measurements} are included. The source instrumentation is recorded in
\href{evidence/signature_transfer/auxfunc-reduced-state.patch}{the
diff}. This changes every SM3 role in the reduced-state experiment and
is not a forgery against the submitted 256-bit profile.

\textbf{{[}E5{]} Relaxed SIS witness.}
\href{code/composite_sis.py}{\emph{composite\_sis.py}},
\href{evidence/composite_sis_5120.json}{run data}, and the
\href{data/sis_witness_5120.npz}{saved witness} implement and record the
full-dimension construction. The witness contains the two matrices
modulo 4 and the nonzero vector; higher matrix bits are unnecessary for
verification. A \href{evidence/composite_sis_replay.json}{publication
replay} reproduced the matrices, vector, and full-modulus lift digest
exactly, excluding timings from the comparison.

\textbf{{[}E6{]} Modular structure.}
\href{code/ring_structure.py}{\emph{ring\_structure.py}} and
\href{evidence/ring_structure.json}{its data} verify the modular
factorization, units, annihilators, Hensel factorization, and one
projection diagnostic.

\textbf{{[}E7{]} Source consistency.}
\href{code/verify_additional_findings.py}{\emph{verify\_additional\_findings.py}}
checks threshold arithmetic and odd-weight challenge candidates. With a
separately obtained submission, its source checker also checks
challenge-stream lengths across nine source trees.
\href{code/verify_original_code.py}{\emph{verify\_original\_code.py}}
invokes that checker and records its results with the C experiments.

\textbf{{[}E8{]} Executed sign-stream discrepancy.}
\href{code/verify_challenge_stream.c}{\emph{verify\_challenge\_stream.c}}
and \href{evidence/original_code.json}{output} give the all-zero-seed
examples in Section 10.1 for all nine implementation/profile
combinations.

\textbf{{[}E9{]} Artifact identities.}
\href{data/submission-sha256.json}{\emph{submission-sha256.json}}
identifies 397 original files: the specification and the implementation
source/build inputs. These third-party inputs are not bundled.
\href{SHA256SUMS}{\emph{SHA256SUMS}} fingerprints the publication
package itself. \href{PROVENANCE.md}{\emph{PROVENANCE.md}} gives
acquisition links and the complete estimator revision.

\textbf{{[}E10{]} Certificate and supplementary checks.}
\href{code/verify_publication_supplement.py}{\emph{verify\_publication\_supplement.py}}
and \href{evidence/verify_publication_supplement.json}{output} verify
the saved SIS witness directly, check the \((n,k)=(2,17)\) index
counterexample, and numerically test the embedding formula on small
instances. Witness verification proves the relation for every
full-modulus lift of the saved modulo-4 matrices.

\textbf{{[}E11{]} Equivalent-secret descriptor.}
\href{code/verify_secret_descriptor.c}{\emph{verify\_secret\_descriptor.c}}
and \href{evidence/original_code.json}{output} check the 141-byte
encoding, the common two-block prefix, and exact reconstruction of every
\(s_1\) expansion byte using a 344-bit descriptor in all six SM3
combinations. The seed is a deterministic test value; no unknown secret
is recovered.

\textbf{{[}E12{]} Separate SIS verifier and small demonstration.}
\href{code/verify_sis_witness.py}{\emph{verify\_sis\_witness.py}} checks
every row, the norm, nonzero status, and response-bound failure of the
saved certificate using NumPy without Sage.
\href{code/sis_mod4_demo.py}{\emph{sis\_mod4\_demo.py}} implements the
two binary solves at small dimensions using only standard Python. Both
include a changed-vector rejection check. Their
\href{evidence/verify_sis_witness.json}{certificate output} and
\href{evidence/sis_mod4_demo.json}{small-example output} are included.

\subsection{12.2 Running the checks}\label{running-the-checks}

From the distribution root, verify package integrity and run the
standard-library checks:

\begin{Shaded}
\begin{Highlighting}[]
\ExtensionTok{python3}\NormalTok{ verify\_package.py}
\ExtensionTok{python3}\NormalTok{ run.py}
\end{Highlighting}
\end{Shaded}

The full mathematical replay additionally requires Sage and NumPy. It
verifies the saved certificate, regenerates the full 5120-dimensional
witness, and compares all mathematical data with the saved records:

\begin{Shaded}
\begin{Highlighting}[]
\ExtensionTok{python3}\NormalTok{ run.py }\AttributeTok{{-}{-}full} \AttributeTok{{-}{-}sage{-}python}\NormalTok{ /path/to/sage{-}env/bin/python}
\end{Highlighting}
\end{Shaded}

Sage and NumPy must be available in that interpreter. Fresh results go
to \emph{verification-output/} by default; \emph{--output-dir} selects
another directory. Recorded publication evidence is preserved.
Wall-clock timings are recorded separately from the comparisons.

To execute {[}E4{]}, {[}E4b{]}, {[}E7{]}, {[}E8{]}, and {[}E11{]}
against the original C implementation, obtain the submission separately
and point to its root:

\begin{Shaded}
\begin{Highlighting}[]
\ExtensionTok{python3}\NormalTok{ run.py }\AttributeTok{{-}{-}submission{-}root}\NormalTok{ /path/to/Octarine}
\end{Highlighting}
\end{Shaded}

The optional signature-transfer model uses about two minutes and up to 3
GB of memory. Include it with \texttt{-\/-signature-transfer}:

\begin{Shaded}
\begin{Highlighting}[]
\ExtensionTok{python3}\NormalTok{ run.py }\AttributeTok{{-}{-}submission{-}root}\NormalTok{ /path/to/Octarine }\DataTypeTok{\textbackslash{}}
  \AttributeTok{{-}{-}signature{-}transfer} \AttributeTok{{-}{-}output{-}dir}\NormalTok{ /tmp/octarine{-}reproduction}
\end{Highlighting}
\end{Shaded}

This verifies all 397 original-file hashes before compilation. It
compiles the small research drivers with GCC and runs nine
challenge-stream configurations, six descriptor configurations, the
SM3-continuation check, and the supplied-state 20-step collision replay
with 64-step controls. The optional 48-bit model builds shipped and
reduced-state variants, validates the complete transfer instance, and
runs full-width and public-artifact controls. It does not run the
submission's build scripts or bundled executables. The source checks can
be combined with the full replay:

\begin{Shaded}
\begin{Highlighting}[]
\ExtensionTok{python3}\NormalTok{ run.py }\AttributeTok{{-}{-}full} \AttributeTok{{-}{-}sage}\NormalTok{ sage }\DataTypeTok{\textbackslash{}}
  \AttributeTok{{-}{-}submission{-}root}\NormalTok{ /path/to/Octarine }\AttributeTok{{-}{-}signature{-}transfer}
\end{Highlighting}
\end{Shaded}

The \href{evidence/summary.json}{reproduction summary} records a
successful run of all these checks. Full environment details, individual
commands, and coverage limits are in
\href{REPRODUCING.md}{\emph{REPRODUCING.md}}. The separate certificate
verifier needs only Python and NumPy:

\begin{Shaded}
\begin{Highlighting}[]
\ExtensionTok{python3}\NormalTok{ code/verify\_sis\_witness.py}
\end{Highlighting}
\end{Shaded}

The artifact checks demonstrate the stated algebra, arithmetic, source
behavior, and saved witness. They do not constitute a full
implementation audit, a full-parameter SM3 collision, recovery of an
unknown signing key, a practical Octarine forgery, or a QROM security
proof.

\subsection{12.3 AI-use disclosure}\label{ai-use-disclosure}

DeepSeek V4 Pro 0831 and Meta Muse Spark 1.3 were used in preparing the
analysis. OpenAI Codex was used for additional technical analysis,
cross-checking of claims, development and execution of reproducibility
code, and preparation of this package. The computational evidence is
available for inspection and replay. Mounir IDRASSI is responsible for
the report and its conclusions. See
\href{AI_DISCLOSURE.md}{\emph{AI\_DISCLOSURE.md}}.

\section*{References}\label{references}
\addcontentsline{toc}{section}{References}

\textbf{{[}S{]}} ARCANE team. \emph{ARCANE-Octarine: Algorithm
Specifications and Supporting Documentation}. Submission snapshot, 53
PDF pages.
\href{https://github.com/ngcc-dev/ngcc-harness/blob/31de7d4d57aca530a65c2c9de6ddaeb25c66548d/sign-16/sign-16-spec.pdf}{Pinned
public copy}. SHA-256 is given in Section 1.1.

\textbf{{[}F202{]}} National Institute of Standards and Technology.
\emph{SHA-3 Standard: Permutation-Based Hash and Extendable-Output
Functions}. FIPS PUB 202, August 2015, Appendix A, Table 4.
\href{https://nvlpubs.nist.gov/nistpubs/FIPS/NIST.FIPS.202.pdf}{Official
publication}.

\textbf{{[}G96{]}} Lov K. Grover. \emph{A Fast Quantum Mechanical
Algorithm for Database Search}. Proceedings of STOC 1996, pp.~212--219.
\href{https://arxiv.org/abs/quant-ph/9605043}{Author preprint}.

\textbf{{[}BHT98{]}} Gilles Brassard, Peter Høyer, and Alain Tapp.
\emph{Quantum Algorithm for the Collision Problem}. LATIN 1998, LNCS
1380, pp.~163--169; preprint 1997.
\href{https://arxiv.org/abs/quant-ph/9705002}{Author preprint}.

\textbf{{[}MNS13{]}} Florian Mendel, Tomislav Nad, and Martin Schläffer.
\emph{Finding Collisions for Round-Reduced SM3}. CT-RSA 2013, LNCS 7779,
pp.~174--188, Section 5.1, Table 3 (p.~182).
\href{https://doi.org/10.1007/978-3-642-36095-4_12}{Publisher record},
\href{https://pure.tugraz.at/ws/portalfiles/portal/80044228/sm3.pdf}{institutional
copy}.

\textbf{{[}BGM16{]}} Andrej Bogdanov, Siyao Guo, Daniel Masny, Silas
Richelson, and Alon Rosen. \emph{On the Hardness of Learning with
Rounding over Small Modulus}. TCC 2016-A, LNCS 9562, pp.~209--224.
\href{https://people.csail.mit.edu/sirichel/lwr.pdf}{Author-hosted
paper}, \href{https://doi.org/10.1007/978-3-662-49096-9_9}{publisher
record}.

\textbf{{[}BRVW26{]}} Joppe W. Bos, Joost Renes, Frederik Vercauteren,
and Peng Wang. \emph{Structured Module Lattice-based Cryptography:
Fine-Grained Parameter Selection without Compromising Performance}. IACR
Cryptology ePrint Archive, Report 2026/098.
\href{https://eprint.iacr.org/2026/098}{Paper record}.

\textbf{{[}B99{]}} Richard P. Brent. \emph{Factorization of the Tenth
Fermat Number}. Mathematics of Computation 68 (1999), pp.~429--451.
\href{https://maths-people.anu.edu.au/~brent/pub/pub161.html}{Author's
publication page}.

\textbf{{[}LE{]}} Martin R. Albrecht et al.~\emph{lattice-estimator}.
Source snapshot identified in the submission by revision \emph{be81b18},
resolved to \emph{be81b183195fb0d62fd882c65620237c6218b981}.
\href{https://github.com/malb/lattice-estimator/blob/be81b183195fb0d62fd882c65620237c6218b981/estimator/sis_lattice.py}{SIS
implementation at that revision}.

\textbf{{[}SP90A{]}} Elaine Barker and John Kelsey. \emph{Recommendation
for Random Number Generation Using Deterministic Random Bit Generators}.
NIST SP 800-90A Revision 1, June 2015.
\href{https://nvlpubs.nist.gov/nistpubs/specialpublications/nist.sp.800-90ar1.pdf}{Official
publication}.

\end{document}
